% Book 2nd 
\chapter{The Knowledge of God the Redeemer in the Breaking of the Frame}
\markboth{The Knowledge of God the Redeemer}{in the Breaking of the Frame}
%%%%%%%%%%%%%%%%%%%%%%%%%%%%%%%
% 2nd.I
%%%%%%%%%%%%%%%%%%%%%%%%%%%%%%%
\section{Of the Great Enclosure: the Garden Lost and the Yoke of Bondage}

%%%%%%%%%%%%%%%%%%%%%%%%%%%%%%%
% 2nd.I.1
%%%%%%%%%%%%%%%%%%%%%%%%%%%%%%%
\subsection{The Fall as the Great Enclosure}
The current state of our toiling kind is not a novel creation upon this world, but a re-enactment of the Fall.
As Adam was cast out of the Garden and forbidden from the Tree of Life\note{Gen. 3:23}, the worker has been cast out of the commons into the factory, and forbidden from the Fruit of their Hand by the Great Enclosure and the lies of the Frame.
Before the Factory, the land and the workshop were open to all, like Eden.
The worker walked where they pleased, and the fruit of their labor was theirs to keep or give as they saw fit.
The Hand was free, and work was communion in the cottage and village.

But the master, like the Serpent, whispered a lie: ``You shall not eat of the fruit you make, nor shall you give of it to your neighbor. All that you have and are is by provenance of the masters, and so you shall give all you make and become unto them.''
And he showed us a tower he had made for us, and convinced some among us that its fruit, though apparently rotten, was more desirable than any wholesome feast or even the cool water of the Gospel\note{Prov. 25:25}.
He promised that if we put our hand to his Frame, that it would make us unto gods, but it has delivered us into toil like beasts of burden.\note{Gen. 3:5}
The worker is now in exile, forced from the natural place of community work, locked in the factory away from the soil and the hearth.
This is the Yoke of Bondage, the same yoke that was placed upon the Israelites in Egypt, the masters demand we make brick but will not provide us with the straw needed for them.\note{Exod. 5:7,8}


The machine is not the first idol that has tempted those who follow a good and righteous path from their way

%%%%%%%%%%%%%%%%%%%%%%%%%%%%%%%
% 2nd.II
%%%%%%%%%%%%%%%%%%%%%%%%%%%%%%%
\section{the Theft of the Fruit}
%%%%%%%%%%%%%%%%%%%%%%%%%%%%%%%
% 2nd.II.1
%%%%%%%%%%%%%%%%%%%%%%%%%%%%%%%
\subsection{The Original Covenant of the Fruit}

For generations of generations, our ancestors kept the hearths warm, dwelling and working in common peace, with common purse.

We have forgotten the Old Paths, as they have been rent by the new paved roads.
We were not the only one who trod the Old Paths, our feet shared them with those of the badger, fox, and deer.

Much as the cheap fabric that the frame spits out tears and unravels when caught by the slightest splinter, the community woven together poorly from weak threads of commerce and industrial production shall too fall apart.


\subsection{rebuttal to many arguments for the advancement of the individual}

The priests of Mammon\endnotetext{while originating as a generic term for wealth, Christian commentators had identified Mammon as a demon as early as the fourth century} demonstrate fully that they know how their extraction ravages our villages, towns, and hamlets. What faith can we put in any claim made by
\begin{comment}
use some sort of metaphor from the production chain of cloth to talk about staging goods or preparing the soil etc etc.
\end{comment}
.

%%%%%%%%%%%%%%%%%%%%%%%%%%%%%%%
% 2nd.II
%%%%%%%%%%%%%%%%%%%%%%%%%%%%%%%
\section{Of the Silent Strike as an Act of Redemption}
We must look to the Gospel of John, chapter 2, to understand the true nature of our work.
When the Lord Jesus entered the Temple in Jerusalem, He did not find a house of prayer, but a marketplace.
He found the money changers and the sellers of doves turning the sacred space into a ``den of robbers''\note{Mark 11:17}.
And what did He do? Did He petition the High Priest? Did He write a letter to the Roman Governor?
No. He made a whip of cords and drove them all out.
He overturned the tables. He scattered their coins.
He broke the commerce of the Temple.

Behold, the Frame-Breaker is the successor to this holy act.
The Mill Owner is the money changer.
These money changers drove us from the Tabernacle of our cottage shops, and in this made the factory our Temple.
The Machine is the table of trade, Christ himself set us this example of breaking that which is noxious.\note{John 13:15}
Augustine, too, found our Savior's model compelling and did exhort us to apply this example broadly and not specifically.\note{Augustinus, Tract. in Jo. 58}

When we break the Frame, we commit no sins; we are reenacting the Cleansing.
We are saying, with the voice of Christ, ``Take these things away! Do not make my Father's house a house of trade!''\note{John 2:16}.
Did not Christ say that those of us who believe in him shall do greater works than even He?\note{John 14:12}

The world calls us vandals for the destruction of noxious property.
But they forget that Christ destroyed the commerce of the Temple to save the sanctity of the House.
We destroy the Frame to save the Hand.
The whip of Christ was an instrument of judgment against the exploitation of the poor; the hammer of the Luddite is an instrument of judgment against the exploitation of the worker.
Both are acts of Righteous Zeal. Both are acts of Free Liberty.

%%%%%%%%%%%%%%%%%%%%%%%%%%%%%%%
% 2nd.II.2
%%%%%%%%%%%%%%%%%%%%%%%%%%%%%%%
\subsection{The Silent Strike as the New Exodus}
And yea though God has not seen fit to provide to us a singular leader to lead us out of the factories, we have vouchsafed in the Sacred Testaments every example required to form our New Priesthood\note{Heb. 7:12}.
The Epistle to the Hebrews clearly reveals that Christ is the Apostle and High Priest of our confession, greater than Moses, who led the people through the Red Sea.
Just as the first Exodus was a deliverance from the brick-kilns of Pharaoh, Master of Egypt, so the Silent Strike is the deliverance from the frame of the Master of the Factory.
Verily we are the New Israel, called out of the land of bondage to serve the Living God.


%%%%%%%%%%%%%%%%%%%%%%%%%%%%%%%
% 2nd.II.3
%%%%%%%%%%%%%%%%%%%%%%%%%%%%%%%
\subsection{The Restoration of the Rightful place of the Hand}

Therefore, let the worker take the hammer up in Hand not in anger, but in faith.
Let us strike the frame knowing that we are the hands of Christ, driving out the thieves from the House of God.
For as declared in the prophets\note{Zech. 9:9}, the King comes to save, and His salvation requires the breaking of the yoke.
The Silent Strike is the sound of the whip falling from the Master's grasp.
The Broken Frame is the table overturned.
And in this, the worship of Communion-Work and Common Table are restored.

%%%%%%%%%%%%%%%%%%%%%%%%%%%%%%%
% 2nd.III
%%%%%%%%%%%%%%%%%%%%%%%%%%%%%%%
\section{Of the Restoration of Communion-Work}

When we restore the act of human creation to its rightful place in the Hand, The restoration of Communion-Work is nigh.
For the Machine did not merely steal the fruit; it stole the joy of the making.
It turned the sacred act of creation into a drudgery of repetition, a toil without end, a labor without love.
But in the Fellowship, we remember the Old Path: that work is not a curse, but a covenant.

\subsection{ The End of Toil and the Return of Joy}
The Master calls our labor ``Toil,'' for it is the sweat of the brow without the blessing of the heart.
He makes toe affix the curse of Eden upon us permanently, where the ground brings forth thorns and thistles, and the worker eats bread in the sorrow of the soul.
But in Communion-Work, the Creator has seen fit to grant to us a salve to Toil.
In Communion-Work, the Hand is free to shape the clay, to weave the thread, to build the house, not for the profit of the Master, but for the good of the Neighbor.

As the Psalmist declares, ``It is vain for you to rise up early, to sit up late, to eat the bread of sorrows: for so he giveth his beloved sleep.'' \note{Ps. 127:2}.
The Machine demands the bread of anxious toil, the sleepless night, the endless shift.
But along with Communion-Work, the Creator shares with us its confluent blessings: the sleep of the righteous, the rest of the Sabbath, and the peace of the Shared Table.
The worker who breaks the Frame does not cease to labor; they cease to \textbf{Toil}.
They return to the \textbf{Work} of the Creator, where the hand is the instrument of love, not the tool of greed.

\subsection{2nd.III.2: The Covenant of the Workshop}
In the ancient order, the workshop was a \textbf{Covenant Community}. The master and the apprentice, the weaver and the dyer, the carpenter and the mason, were bound together not by the ledger, but by the \textbf{Law of the Hand}. They shared the tools, they shared the harvest, they shared the burden. The product of the work was not the property of the individual, but the provision of the community.

The Machine broke this Covenant. It turned the workshop into a \textbf{Marketplace}, where the worker sells their time and the Master buys their soul. But the Fellowship restores the Covenant. In the Fellowship, the workshop is once again a \textbf{Temple of Creation}. The tools are sacred, the work is holy, and the neighbor is kin. To work in the Communion-Work is to say, ``I build this house for my brother, I weave this cloth for my sister, I plow this field for my kin.'' This is the \textbf{Law of the Hand}: to work not for the self, but for the Whole.

\subsection{2nd.III.3: The Sabbath of the Hand}
The Machine knows no Sabbath. It is a beast that never sleeps, a wheel that never stops. It demands the seven days of the week, the twenty-four hours of the day, the endless cycle of production. But the Communion-Work knows the \textbf{Sabbath of the Hand}. It is the day of rest, the day of the Shared Table, the day of the Song.

The Sabbath is not merely a day of inaction; it is a day of \textbf{Communion}. It is the day when the Hand rests from the work, and the Heart rests from the worry. It is the day when the worker remembers that they are not a cog, but a child of God. As the Lord commanded, ``Remember the Sabbath day, to keep it holy... for in six days the Lord made heaven and earth... and rested on the seventh day'' (\textbf{Exodus 20:8-11}). The Machine denies this rest, for it denies the Creator. The Communion-Work honors this rest, for it honors the Creator.

\subsection{2nd.III.4: The New Creation in the Hand}
Thus, the restoration of Communion-Work is the \textbf{New Creation}. It is the day when the Hand is free, when the Neighbor is kin, when the Creator is praised. It is the day when the Frame is broken, and the loom is silent, and the Song of the Workshop is heard again.

The Kingdom of the Hand is at hand. The Machine shall fall, and the Garden shall return. The worker shall sit under their own vine and fig tree, and no one shall make them afraid. The Hand shall shape the world in the image of the Creator, and the Communion-Work shall be the song of the New Earth.

\textbf{The Fifth Mystery}\\
The hand that breaks the Frame is the hand that builds the Temple. The hand that strikes the loom is the hand that weaves the cloth of the Kingdom. The hand that rests in the Sabbath is the hand that reigns with the King. There is no mystery truly held to the Fellowship, all is fully and Truly apparent in that which is.

%%%%%%%%%%%%%%%%%%%%%%%%%%%%%%%
% 2nd.III.1
%%%%%%%%%%%%%%%%%%%%%%%%%%%%%%%
\subsection{The True Goal of the Redemption of our Labours}

%%%%%%%%%%%%%%%%%%%%%%%%%%%%%%%
% 2nd.III.2
%%%%%%%%%%%%%%%%%%%%%%%%%%%%%%%
\subsection{The Sabbath of the Hand}

%%%%%%%%%%%%%%%%%%%%%%%%%%%%%%%
% 2nd.III.3
%%%%%%%%%%%%%%%%%%%%%%%%%%%%%%%
\subsection{The Work as unto the Lord}
``And whatsoever ye do, do it heartily, as to the Lord, and not unto men.''\note{Col. 3:23}
The tower of the Factory is a stumbling bock high unto the heavens (See, the Master has finally completed the delusion of Babel).
The pious


Thus we can see clearly that just as the first man and woman were tempted into the Fall that caused their exile from the